\documentclass[12pt]{article}

% Required packages
\usepackage[utf8]{inputenc}
\usepackage[margin=1in]{geometry}
\usepackage{hyperref}
\usepackage{graphicx}
\usepackage{longtable}
\usepackage{booktabs}
\usepackage{array}
\usepackage{enumitem}
\usepackage{xcolor}

% Define colors for comments/instructions
\definecolor{instructionblue}{RGB}{0,0,255}
\newcommand{\instruction}[1]{\textcolor{instructionblue}{\emph{#1}}}

% Hyperref setup
\hypersetup{
    colorlinks=true,
    linkcolor=blue,
    filecolor=magenta,
    urlcolor=cyan,
    citecolor=blue,
}

\begin{document}

% ===================================================================
% ARTICLE INFORMATION
% ===================================================================

\title{%
% INSTRUCTION: The article title MUST include the word 'data' or 'dataset'
Peru GDP Real-Time Dataset (1994-2025): Tracking Three Decades of Revisions
}

\author{%
Jason Cruz\textsuperscript{1,*}, Diego Winkelried\textsuperscript{1}, Javier Torres\textsuperscript{1} \\[1mm]
\small\textsuperscript{1}Centro de Investigación (CIUP), Universidad del Pacífico, \\
\small Av. Salaverry 2020, Jesús María, Lima 15072, Peru \\[2mm]
\small\textsuperscript{*}Corresponding author \\
\small Email: jj.cruza@up.edu.pe
}

\date{} % Leave empty for journal to handle

\maketitle

% ===================================================================
% KEYWORDS
% ===================================================================
\noindent\textbf{Keywords:}
% INSTRUCTION: Include 4-8 keywords, separated by semi-colons. Do not repeat words from title.
Real-time macroeconomic statistics; GDP revisions; Emerging markets; Nowcasting; Vintage data; Data curation; Central bank publications; Economic indicators

% ===================================================================
% ABSTRACT
% ===================================================================
\begin{abstract}
% INSTRUCTION: 100-500 words. Describe data collection process, dataset, and reuse potential.
% Do NOT provide conclusions or interpretations.

This data article presents a comprehensive real-time dataset (RTD) of Peru's Gross Domestic Product (GDP) growth rates spanning 1994 to 2025, comprising over 1000 data releases across 30+ years. The dataset was constructed by systematically collecting and processing monthly GDP tables from the Central Reserve Bank of Peru (BCRP) Weekly Reports. For the 1994-2012 period, data were extracted from scanned hardcover volumes using Optical Character Recognition (OCR) techniques with extensive manual verification. From 2013 onwards, data were collected via automated web scraping and PDF table extraction from digital publications. The dataset includes monthly, quarterly, and annual GDP growth rates for aggregate GDP and 8 economic sectors, organized in two complementary formats: vintage format (columns = release dates) and releases format (columns = revision sequences). Additional variants account for base-year adjustments and identify benchmark revisions. All data are provided in CSV format with comprehensive documentation. An open-source Python pipeline enables full reproducibility and future updates. This dataset enables research on GDP revision patterns, real-time forecasting accuracy, nowcasting model evaluation, and cross-country comparisons of statistical practices in emerging economies.

\end{abstract}

\newpage

% ===================================================================
% SPECIFICATIONS TABLE
% ===================================================================
\section*{Specifications Table}

% INSTRUCTION: This table provides structured metadata about your dataset

\begin{longtable}{|>{\raggedright\arraybackslash}p{3.5cm}|>{\raggedright\arraybackslash}p{11cm}|}
\hline
\textbf{Subject} & Economics \\
\hline
\textbf{Specific subject area} &
% INSTRUCTION: Max 150 characters without spaces
Real-time macroeconomic data, GDP revisions, emerging market economic statistics \\
\hline
\textbf{Type of data} &
% INSTRUCTION: List data types and formats
Table, Raw, Analyzed, Filtered \\
CSV files \\
\hline
\textbf{Data collection} &
% INSTRUCTION: Describe instruments and methods. Max 600 characters without spaces.
Data collected from BCRP Weekly Reports (1994-2025). Pre-2013: scanned PDFs processed with Tesseract OCR, manual verification, 70+ cleaning functions. Post-2013: Selenium web scraping, Tabula-py PDF extraction, automated cleaning pipeline. Monthly GDP tables systematically extracted, standardized, and organized by release date. Python-based processing with quality validation checks. \\
\hline
\textbf{Data source location} &
% INSTRUCTION: Where data were collected or stored
Institution: Universidad del Pacífico - CIUP \\
City/Country: Lima, Peru \\
Geographic coordinates: -12.0897°S, -77.0428°W \\
Data source: Banco Central de Reserva del Perú (BCRP) \\
URL: https://www.bcrp.gob.pe/publicaciones/nota-semanal.html \\
\hline
\textbf{Data accessibility} &
% INSTRUCTION: All data must be in a public repository BEFORE publication
Repository name: Zenodo \\
Data identification number: 10.5281/zenodo.18099975 \\
Direct URL to data: https://doi.org/10.5281/zenodo.18099975 \\
Related code repository: https://github.com/JasonCruz18/peru\_gdp\_revisions \\
Instructions: Dataset freely accessible via DOI. 16 CSV files with README documentation. \\
\hline
\textbf{Related research article} &
% INSTRUCTION: Cite ONE related article. Upload a copy with submission. If none, state 'none'.
None \\
\hline
\end{longtable}

% ===================================================================
% VALUE OF THE DATA
% ===================================================================
\section*{Value of the Data}

% INSTRUCTION: Provide 3-6 bullet points. Max 150 words each. No conclusions/inferences.
% Answer: Why valuable? How can it be reused?

\begin{itemize}[leftmargin=*]

\item \textbf{First real-time GDP dataset for Peru and one of the few for emerging markets.} Existing real-time datasets focus primarily on advanced economies (US, Europe). This dataset fills a critical gap by providing systematic vintage data for an emerging market economy, enabling comparative studies of revision patterns across development levels and statistical capacity.

\item \textbf{Enables evaluation of real-time forecasting and nowcasting models.} Researchers can assess forecast accuracy using information available at the time predictions were made, rather than revised data. This is essential for developing robust nowcasting models for emerging economies where data quality and revision patterns differ substantially from advanced economies.

\item \textbf{Supports revision analysis and statistical quality assessment.} The dataset allows investigation of revision properties (bias, efficiency, rationality), identification of benchmark revisions, and analysis of base-year adjustments. This informs best practices for data producers and helps researchers understand data generation processes.

\item \textbf{Fully reproducible and updatable.} Complete open-source Python pipeline enables exact replication and seamless extension as new data become available. Modular architecture supports adaptation to other countries or statistical agencies with minimal modifications.

\item \textbf{Multiple formats optimized for different research applications.} Vintage format facilitates real-time analysis; releases format enables revision econometrics; benchmark indicators isolate methodological changes; base-year adjustments handle structural breaks. Researchers can select the appropriate variant for their specific analytical needs.

\item \textbf{Comprehensive documentation and quality assurance.} Extensive metadata, cleaning procedures, and validation checks are fully documented. Dual OCR/automated extraction methods ensure coverage across three decades despite changing data formats. Dataset includes 1000+ vintages across 8 economic sectors.

\end{itemize}

% ===================================================================
% BACKGROUND
% ===================================================================
\section*{Background}

% INSTRUCTION: Max 200 words. Describe motivation and context. No conclusions/interpretations.
% If related to a research article, describe how this data article adds value.

Real-time datasets preserve the actual information available to policymakers and forecasters at specific points in time, enabling accurate evaluation of forecast performance and analysis of data revision patterns. While such datasets are well-established for advanced economies, they remain scarce for emerging markets. Peru provides an ideal case study: it has maintained consistent GDP reporting through the Central Reserve Bank's Weekly Reports for over 30 years, yet no systematic real-time dataset existed.

This dataset was compiled to address this gap and support research on GDP revisions, nowcasting, and forecast evaluation in emerging economies. The construction process required developing specialized methods to handle heterogeneous data formats: scanned documents (1994-2012) necessitating OCR techniques, and digital PDFs (2013-2025) amenable to automated extraction. The resulting dataset provides complete coverage of monthly, quarterly, and annual GDP growth rates across all major economic sectors, organized in formats suitable for both real-time forecasting applications and econometric analysis of revision properties. The accompanying open-source pipeline ensures reproducibility and facilitates future updates and extensions to other emerging market economies.

% ===================================================================
% DATA DESCRIPTION
% ===================================================================
\section*{Data Description}

% INSTRUCTION: Describe your dataset structure. Refer to all folders/subfolders/files individually.
% Use visual aids (tables, graphs, figures with captions). No interpretations/conclusions.

The dataset comprises 16 CSV files organized in two main formats (vintage and releases), each available in four variants to support different analytical needs. All files are hosted on Zenodo (DOI: 10.5281/zenodo.18099975) with accompanying README documentation.

\subsection*{File Organization}

Table \ref{tab:file_structure} presents the complete file structure. Files are named following a consistent convention: \texttt{[frequency]\_gdp\_[format]\_[modifier].csv}, where frequency is monthly or quarterly, format is vintages or releases, and optional modifiers indicate adjusted (base-year adjustments) or benchmark (methodological revision indicators) variants.

\begin{table}[h]
\centering
\caption{Dataset file structure}
\label{tab:file_structure}
\begin{tabular}{lrrr}
\toprule
File Name & Rows & Columns & Size (MB) \\
\midrule
\multicolumn{4}{l}{\textbf{Monthly Data}} \\
monthly\_gdp\_vintages.csv & 1000+ & 300+ & 3.0 \\
monthly\_gdp\_vintages\_adjusted.csv & 1000+ & 300+ & 3.0 \\
monthly\_gdp\_vintages\_benchmark.csv & 1000+ & 300+ & 3.0 \\
monthly\_gdp\_releases.csv & 300+ & 72 & 0.4 \\
monthly\_gdp\_releases\_adjusted.csv & 300+ & 72 & 0.5 \\
monthly\_gdp\_releases\_benchmark.csv & 300+ & 72 & 0.4 \\
\midrule
\multicolumn{4}{l}{\textbf{Quarterly \& Annual Data}} \\
quarterly\_gdp\_vintages.csv & 1000+ & 200+ & 2.5 \\
quarterly\_gdp\_vintages\_adjusted.csv & 1000+ & 200+ & 2.5 \\
quarterly\_gdp\_vintages\_benchmark.csv & 1000+ & 200+ & 2.5 \\
quarterly\_gdp\_releases.csv & 200+ & 60 & 0.4 \\
quarterly\_gdp\_releases\_adjusted.csv & 200+ & 60 & 0.4 \\
quarterly\_gdp\_releases\_benchmark.csv & 200+ & 60 & 0.4 \\
\bottomrule
\end{tabular}
\end{table}

\subsection*{Data Formats}

\textbf{Vintage Format:} Each row represents an economic sector; columns represent release dates (vintages). Cell values are GDP growth rates as published in that vintage. This format preserves the exact information available at any historical point in time, essential for real-time forecasting applications.

\textbf{Releases Format:} Each row represents a target period (reference month/quarter); columns represent revision sequences (1st release, 2nd release, etc.). This format facilitates analysis of revision patterns and computation of revision statistics.

\subsection*{Data Variants}

\textbf{Base Variant:} Standard GDP growth rates as published, suitable for most applications.

\textbf{Adjusted Variant:} Applies sentinel value (-999999.0) to growth rates that become incomparable due to base-year changes (1994, 2007, 2019). Marks invalid observations without removing them, preserving complete revision history while flagging structural breaks.

\textbf{Benchmark Variant:} Replaces growth rates with binary indicators (1.0 = benchmark revision with simultaneous update of monthly and quarterly/annual tables; 0.0 = routine update). Enables isolation of methodological revisions from standard data updates.

\subsection*{Coverage}

\begin{itemize}
\item \textbf{Temporal coverage:} 1994-2025 (31 years)
\item \textbf{Frequencies:} Monthly, quarterly, annual
\item \textbf{Vintages tracked:} 1000+ data releases
\item \textbf{Economic sectors:} 9 total (aggregate GDP + 8 sectors)
\begin{itemize}
\item Agriculture and livestock
\item Fishing
\item Mining and fuel
\item Manufacturing
\item Electricity and water
\item Construction
\item Commerce
\item Other services
\end{itemize}
\end{itemize}

\subsection*{Data Schema}

\textbf{Vintage Format Schema:}
\begin{itemize}
\item \texttt{industry} (string): Sector name (standardized)
\item \texttt{vintage} (string): Publication month in YYYYmM format (e.g., 2020m5)
\item \texttt{tp\_YYYYmM} / \texttt{tp\_YYYYqN} / \texttt{tp\_YYYY} (float): Growth rate columns for each target period
\end{itemize}

\textbf{Releases Format Schema:}
\begin{itemize}
\item \texttt{target\_period} (string): Reference period in YYYYmM/YYYYqN/YYYY format
\item \texttt{industry\_h} (float): $h$-th release for each industry (e.g., \texttt{agriculture\_1}, \texttt{agriculture\_2}, etc.)
\end{itemize}

\subsection*{Missing Values}

Missing values (NaN) occur legitimately in vintage format when:
\begin{itemize}
\item Data not yet published for future periods (normal in real-time datasets)
\item Historical data outside the reporting window in early vintages
\item Sectoral data not available for certain periods
\end{itemize}

Sentinel value (-999999.0) in adjusted variants marks observations affected by base-year changes, not true missing values.

% ===================================================================
% EXPERIMENTAL DESIGN, MATERIALS AND METHODS
% ===================================================================
\section*{Experimental Design, Materials and Methods}

% INSTRUCTION: Comprehensive description of data acquisition. No character limit.
% Include all code, software, tools, instruments, conditions. Use figures/tables.
% No interpretations/conclusions.

\subsection*{Data Sources}

\textbf{Primary Source:} Banco Central de Reserva del Perú (BCRP) Weekly Reports (Nota Semanal), published every Friday since 1994. Each report reproduces GDP growth statistics from Peru's National Institute of Statistics and Informatics (INEI) in standardized tables.

\textbf{Historical Context:} Prior to 2013, Weekly Reports existed only as hardcover volumes stored at the BCRP's Renzo Rossini Library (Lima, Peru). From 2013 onwards, reports became freely available as digital PDFs at \url{https://www.bcrp.gob.pe/publicaciones/nota-semanal.html}.

\subsection*{Revision Calendar}

Peru lacks an official GDP revision calendar. We reconstructed an implicit monthly revision calendar by systematically collecting the last Weekly Report of each month, which defines a data vintage. This approach is justified by: (1) INEI Chief Resolution No. 316-2003-INEI mandating quarterly updates, and (2) empirical observation that Weekly Reports update GDP at least monthly.

\subsection*{Data Collection Methods}

Data collection employed two distinct methodologies depending on time period:

\subsubsection*{Pre-2013 Data (Scanned PDFs - OCR-based Extraction)}

\textbf{Physical Access:} Hardcover Weekly Report volumes were consulted in-person at the BCRP Renzo Rossini Library. Volumes were photocopied or photographed (no removal permitted).

\textbf{OCR Processing:} Scanned PDFs were processed using Tesseract OCR engine (version 4.1+, maintained by Google) with extensive preprocessing:
\begin{enumerate}
\item PNG conversion at minimum 300 DPI resolution
\item Grayscale conversion and adaptive binarization
\item Noise removal via median filtering
\item Skew correction (deskewing) for rotated text
\item Contrast enhancement using CLAHE algorithm
\item Border removal to prevent false detections
\item Character thinning and skeletonization
\end{enumerate}

\textbf{Quality Challenges:} Scanned sources exhibited warped text, misaligned tables, shadowing, blur, highlighted entries, inconsistent fonts, and residual marks obscuring digits/decimal points. These issues necessitated extensive manual verification.

\textbf{Manual Verification:} Each OCR-extracted table was manually reviewed and corrected for misclassifications (e.g., 9.4 read as 5.4). Corrected data were stored in CSV format.

\textbf{Automated Cleaning:} Pre-2013 tables were processed using Python module \texttt{peru\_gdp\_rtd.cleaners.old\_table\_cleaner} with orchestration functions \texttt{old\_table\_1\_runner} and \texttt{old\_table\_2\_runner}.

\subsubsection*{Post-2013 Data (Digital PDFs - Automated Pipeline)}

\textbf{Web Scraping:} Automated download of PDF reports using Selenium WebDriver (version 4.0+) with Chrome browser. Scraper implements rate limiting, retry logic with exponential backoff, and concurrent downloads respecting server policies.

\textbf{PDF Table Extraction:} Tables extracted using Tabula-py (version 2.5+), a Python wrapper for Tabula Java library. Keyword-based page filtering identifies relevant GDP tables before extraction.

\textbf{Common Extraction Issues:} Digital PDFs presented encoded characters (e.g., \texttt{(cid:101)}), incomplete table borders, duplicate documents, empty cells misinterpreted as headers, missing year labels, and fragmented borders.

\textbf{Automated Cleaning:} Post-2013 tables processed using Python module \texttt{peru\_gdp\_rtd.cleaners.new\_table\_cleaner} with orchestration functions \texttt{new\_table\_1\_runner} and \texttt{new\_table\_2\_runner}.

\subsection*{Data Processing Pipeline}

The complete pipeline is implemented in Python 3.10+ and organized into 14 specialized modules totaling 4000+ lines of documented code. All processing is configuration-driven via YAML file (\texttt{config/config.yaml}), eliminating hardcoding.

\textbf{Pipeline Architecture:}
\begin{enumerate}
\item \textbf{Configuration Module} (\texttt{peru\_gdp\_rtd.config}): Type-safe settings management with Pydantic-style validation

\item \textbf{Scrapers Module} (\texttt{peru\_gdp\_rtd.scrapers}): Selenium-based PDF downloading with authentication, rate limiting, and retry logic

\item \textbf{Processors Module} (\texttt{peru\_gdp\_rtd.processors}): PDF processing, keyword extraction, metadata parsing

\item \textbf{Cleaners Module} (\texttt{peru\_gdp\_rtd.cleaners}): 70+ specialized cleaning functions organized in 7 sub-modules:
\begin{itemize}
\item \texttt{text\_cleaners} (4 functions): Unicode normalization, whitespace handling
\item \texttt{table\_cleaners} (22 functions): DataFrame operations, noise removal
\item \texttt{column\_handlers} (14 functions): Column renaming, type conversion
\item \texttt{table1\_cleaners} (13 functions): Monthly table transformations
\item \texttt{table2\_cleaners} (13 functions): Quarterly/annual table transformations
\item \texttt{old\_table\_cleaner}: Orchestrator for pre-2013 data
\item \texttt{new\_table\_cleaner}: Orchestrator for post-2013 data
\end{itemize}

\item \textbf{Transformers Module} (\texttt{peru\_gdp\_rtd.transformers}):
\begin{itemize}
\item \texttt{vintage\_preparator}: Reshapes tables to vintage format
\item \texttt{concatenator}: Merges vintages across years (372 lines)
\item \texttt{metadata\_handler}: Base-year mapping, benchmark indicators (655 lines)
\item \texttt{releases\_converter}: Vintage-to-releases transformation (241 lines)
\end{itemize}

\item \textbf{Orchestration Module} (\texttt{peru\_gdp\_rtd.orchestration}): High-level workflow coordination

\item \textbf{Utilities Module} (\texttt{peru\_gdp\_rtd.utils}): RecordManager for idempotent operations, alerts
\end{enumerate}

\textbf{Execution:} Complete pipeline executed via single command: \texttt{python scripts/update\_rtd.py}. Supports selective step execution, dry-run mode, and custom configuration files.

\subsection*{Data Cleaning Procedures}

\textbf{Text Normalization:}
\begin{itemize}
\item Unicode normalization (NFC form)
\item Whitespace standardization (tabs → spaces)
\item Special character handling
\item Case normalization for sector names
\end{itemize}

\textbf{Table Structure Cleaning:}
\begin{itemize}
\item Empty row removal (NA or whitespace-only)
\item Duplicate detection and removal
\item Header row extraction
\item Column alignment correction
\item Numeric column identification
\end{itemize}

\textbf{Numerical Data Cleaning:}
\begin{itemize}
\item Encoded character replacement
\item Decimal separator normalization (comma → period)
\item Thousands separator removal
\item String-to-float conversion with error handling
\item Missing value handling (empty strings, dashes → NaN)
\end{itemize}

\textbf{Temporal Data Handling:}
\begin{itemize}
\item Missing year label inference via forward-filling
\item Spanish month abbreviation parsing (ene, feb, mar...)
\item Quarter identification (I, II, III, IV)
\item Target period construction (e.g., 2020m5, 2020q2)
\end{itemize}

\textbf{Sector Name Standardization:} Comprehensive mappings handle Spanish/English terminology variations across vintages (e.g., ``agropecuario'' → ``agriculture'', ``minería e hidrocarburos'' → ``mining'').

\subsection*{Base-Year Adjustments}

Peru's GDP series underwent three base-year changes (1994, 2007, 2019), creating structural breaks. We identify affected periods via metadata mapping and apply sentinel value (-999999.0) to mark incomparable observations. This preserves complete revision history while flagging invalid comparisons.

\subsection*{Benchmark Indicators}

Benchmark revisions occur when INEI simultaneously updates monthly and quarterly/annual tables, typically due to methodological changes. We identify these events by comparing release dates across table types and generate binary indicators (1.0 = benchmark, 0.0 = routine update).

\subsection*{Format Conversion}

Vintage-to-releases conversion implemented via NumPy-based vertical alignment algorithm. For each target period, the algorithm extracts all releases chronologically and aligns them in revision sequence columns.

\subsection*{Quality Validation}

Automated validation checks include:
\begin{itemize}
\item Schema validation (expected columns, correct types)
\item Range validation (growth rates within plausible bounds $\pm$50\%)
\item Continuity checks (vintage date monotonicity)
\item Completeness checks (missing value patterns)
\end{itemize}

\subsection*{Software and Tools}

\textbf{Core Dependencies:}
\begin{itemize}
\item Python 3.10+ (programming language)
\item pandas 2.0+ (DataFrame operations)
\item numpy 1.24+ (numerical computations)
\item selenium 4.0+ (web scraping)
\item tabula-py 2.5+ (PDF table extraction)
\item tesseract-ocr 4.1+ (OCR engine for pre-2013 data)
\item PyMuPDF 1.23+ (PDF metadata extraction)
\item tqdm 4.65+ (progress tracking)
\item pyyaml 6.0+ (configuration parsing)
\end{itemize}

\textbf{Development Environment:}
\begin{itemize}
\item Operating System: Windows 10, macOS, Linux (cross-platform tested)
\item Java Runtime Environment 8+ (required for Tabula)
\item Git for version control
\item GitHub Actions for continuous integration
\end{itemize}

\textbf{Hardware Requirements:}
\begin{itemize}
\item CPU: Multi-core processor recommended
\item RAM: 8 GB recommended
\item Disk: 5 GB for full dataset archive
\item Internet: Stable connection for PDF downloads
\end{itemize}

\subsection*{Reproducibility}

Complete replication ensured through:
\begin{itemize}
\item Open-source code repository: \url{https://github.com/JasonCruz18/peru_gdp_revisions}
\item MIT License (permits commercial and academic use)
\item Pinned dependency versions (\texttt{requirements-frozen.txt})
\item Configuration-driven execution (no hardcoding)
\item Comprehensive documentation (130+ pages)
\item Tutorial Jupyter notebooks (6 step-by-step guides)
\item Continuous integration testing (cross-platform, Python 3.10-3.12)
\item Machine-readable citation metadata (CITATION.cff)
\end{itemize}

Estimated runtime: 45-60 minutes for complete pipeline; 5-10 minutes for incremental updates.

% ===================================================================
% LIMITATIONS
% ===================================================================
\section*{Limitations}

% INSTRUCTION: Max 200 words. Describe data collection/curation problems, quality issues, biases.
% No analysis/interpretation limitations. If none, write 'None' or 'Not applicable'.

This dataset has several limitations stemming from source data characteristics and collection constraints:

\textbf{OCR Accuracy:} Pre-2013 data extraction via OCR is subject to recognition errors despite extensive preprocessing and manual verification. Poor scan quality (warping, shadowing, blur) in some hardcover volumes may introduce residual errors in early vintages.

\textbf{Publication Irregularities:} While we reconstructed a monthly revision calendar, actual Weekly Report publication was sometimes irregular, particularly in early years. Missing or incomplete reports may create gaps in vintage coverage.

\textbf{Sectoral Coverage:} Sector definitions and naming conventions changed over time, requiring manual mapping. Some sectors were not reported consistently across all vintages, resulting in missing values for certain sector-period-vintage combinations.

\textbf{Base-Year Changes:} Three base-year revisions (1994, 2007, 2019) create structural breaks making cross-period comparisons invalid. While we flag affected observations, users must carefully handle these discontinuities in analysis.

\textbf{Metadata Completeness:} Revision metadata (reasons for revisions, methodological changes) are not systematically documented by the source and therefore not included in this dataset. Benchmark indicators identify simultaneous table updates but cannot distinguish specific methodological changes.

\textbf{Representativeness:} This dataset reflects only officially published GDP statistics. Unpublished revisions or corrections are not captured.

% ===================================================================
% ETHICS STATEMENT
% ===================================================================
\section*{Ethics Statement}

% INSTRUCTION: Select relevant statement(s). Delete those not applicable.

The authors confirm that this work does not involve human subjects, animal experiments, or data collected from social media platforms. All data are derived from publicly available official government publications (BCRP Weekly Reports), which reproduce statistics from Peru's National Institute of Statistics and Informatics (INEI). No ethical approval was required for this data compilation work. The authors have read and followed the ethical requirements for publication in Data in Brief.

% ===================================================================
% CRediT AUTHOR STATEMENT
% ===================================================================
\section*{CRediT Author Statement}

% INSTRUCTION: Outline contributions using CRediT categories
% See: https://www.elsevier.com/authors/journal-authors/policies-and-ethics/credit-author-statement

\textbf{Jason Cruz:} Conceptualization, Data curation, Formal analysis, Investigation, Methodology, Software, Validation, Visualization, Writing - original draft, Writing - review \& editing.

\textbf{Diego Winkelried:} Conceptualization, Formal analysis, Methodology, Supervision, Writing - review \& editing.

\textbf{Javier Torres:} Conceptualization, Formal analysis, Methodology, Supervision, Writing - review \& editing.

% ===================================================================
% ACKNOWLEDGEMENTS
% ===================================================================
\section*{Acknowledgements}

% INSTRUCTION: Mention contributors who don't meet authorship criteria.
% List funding sources in standard format.

This research was funded by the Consejo Nacional de Ciencia, Tecnología e Innovación Tecnológica (CONCYTEC) through PROCIENCIA under grant E041-2025-04 (Proyectos de Investigación en Ciencias Sociales, contract PE501096145-2025).

We thank the staff of the BCRP Renzo Rossini Library for facilitating access to historical Weekly Report volumes. We also acknowledge the BCRP for maintaining publicly accessible digital archives of recent publications.

% ===================================================================
% DECLARATION OF COMPETING INTERESTS
% ===================================================================
\section*{Declaration of Competing Interests}

% INSTRUCTION: Choose one statement and delete the other

The authors declare that they have no known competing financial interests or personal relationships that could have appeared to influence the work reported in this paper.

% ===================================================================
% REFERENCES
% ===================================================================
\section*{References}

% INSTRUCTION: Maximum 20 references. Format as numbered list.
% In-text citations: number(s) in square brackets [1], [2-4], etc.
% If supporting a research article, cite it first.
% MUST cite your dataset in the repository.

\begin{enumerate}

\item Cruz J, Winkelried D, Torres J. Peru GDP Real-Time Dataset (1994-2025). Zenodo; 2025. DOI: 10.5281/zenodo.18099975. Available from: \url{https://doi.org/10.5281/zenodo.18099975}

\item Croushore D, Stark T. A Real-Time Data Set for Macroeconomists. Journal of Econometrics. 2001;105(1):111-130. DOI: 10.1016/S0304-4076(01)00072-0

\item Croushore D. Frontiers of Real-Time Data Analysis. Journal of Economic Literature. 2011;49(1):72-100. DOI: 10.1257/jel.49.1.72

\item Banco Central de Reserva del Perú. Weekly Reports (Nota Semanal). Available from: \url{https://www.bcrp.gob.pe/publicaciones/nota-semanal.html} [Accessed 2025-12-31]

\item Instituto Nacional de Estadística e Informática (INEI). National Accounts Methodology. Available from: \url{https://www.inei.gob.pe/estadisticas/indice-tematico/economia/} [Accessed 2025-12-31]

\end{enumerate}

\end{document}
